%%%%%%%%%%%%%%%%%%%%%%%%%%%%%%%%%%%%%%%%%%%%%%%%%%%%%%%%%%%%%%%%%%%%%%%%%%%%%%
%  Unit III appendices
%%%%%%%%%%%%%%%%%%%%%%%%%%%%%%%%%%%%%%%%%%%%%%%%%%%%%%%%%%%%%%%%%%%%%%%%%%%%%%

\startappendices

\chapter{Syllabus-to-Chapter Concordance}

This appendix exists so that coverage can be \emph{audited} rather than assumed.
The left-hand column reproduces the Unit~III syllabus descriptor phrase by
phrase, in the order the official document prints it. The right-hand column gives
the chapter and section that carries it. Every phrase resolves to a location, and
no numbered chapter of this volume contains material that is not traceable to a
phrase in this table.

\begin{mmlongtable}{Concordance for Unit III}{tab:concordance3}%
{@{}P{0.36\textwidth}P{0.57\textwidth}@{}}
\rowcolor{ocre}\thd{Syllabus phrase} & \thd{Where it is covered}\\
\midrule
\endfirsthead
\rowcolor{ocre}\thd{Syllabus phrase} & \thd{Where it is covered}\\
\midrule
\endhead
\multicolumn{2}{@{}l}{\bfseries\color{ocredark} First half --- Web Analytics \& Google Analytics}\\
\midrule
\rlab{Web Analytics} \newline {\scriptsize the heading concept} & Chapter~1
entire. Definition and importance \S\,1.1; traditional against web-analytics
based analysis \S\,1.2; the three metric families \S\,1.3; core metrics
\S\,1.4; bounce rate \S\,1.5; the dual lens \S\,1.6; the role of web analytics
\S\,1.7; benefits \S\,1.8; KPI alignment \S\,1.9; A/B testing and CRO
\S\,1.10.\\
\rowcolor{tablealt}
\rlab{Google Analytics Tools} & Chapter~2 entire. Purpose \S\,2.1; key features
\S\,2.2; acquisition channels \S\,2.3; report families \S\,2.4; setting up and
tracking goals \S\,2.5; Google Tag Manager \S\,2.6.\\
\rlab{Web Analytics Tools} & Chapter~3 entire. Tools other than Google Analytics
\S\,3.1; combining quantitative and qualitative tools \S\,3.2.\\
\midrule
\multicolumn{2}{@{}l}{\bfseries\color{ocredark} Second half --- E-mail Marketing}\\
\midrule
\rlab{Effective E-mail Campaigns} & Chapter~4 entire. Definition \S\,4.1; E-mail
1.0 against 2.0 \S\,4.2; advantages \S\,4.3; disadvantages \S\,4.4; opt-in
against spam \S\,4.5; anatomy of an effective e-mail \S\,4.6; the nine-type
content arsenal \S\,4.7; transactional against promotional \S\,4.8.\\
\rowcolor{tablealt}
\rlab{E-mail Plan} & Chapter~5 entire. The seven nodes \S\,5.1; a worked plan
\S\,5.2; segmentation and personalisation \S\,5.3.\\
\rlab{E-mail Marketing Campaign Analysis} & Chapter~6 entire. What campaign
analysis means \S\,6.1; the measurement funnel \S\,6.2; interpreting the KPIs
\S\,6.3; challenges and solutions \S\,6.4; raising open and click rates
\S\,6.5.\\
\rowcolor{tablealt}
\rlab{Measuring Conversions \& keeping up} & Chapter~7 entire. Offline tracking
\S\,7.1; integrating analytics with e-mail \S\,7.2; the technology stack
\S\,7.3; responsive design \S\,7.4; the four pillars \S\,7.5.\\
\end{mmlongtable}

\begin{keybox}[No Off-Syllabus Appendix in This Volume]
Unlike Volumes~1 and~2, every topic in the source material for Unit~III maps to a
phrase in the unit descriptor. Nothing had to be set aside, so this volume has no
appendix of supplementary examined topics.
\end{keybox}

\chapter{Previous-Year Question Bank --- Unit III}

Every question in this appendix is reproduced from an actual RVCE Marketing
Management Continuous Internal Evaluation or Semester End Examination paper for
this course. Answers and pointers follow the official schemes of valuation where
those were released.

\section*{Part A --- Two-Mark Questions}
\addcontentsline{toc}{section}{Part A --- Two-Mark Questions}

\begin{enumerate}[leftmargin=2.2em,itemsep=1.1ex]

\item \textbf{Define Web Analytics and mention its importance in digital
      marketing activities.} \pyqr{CIE-II, 26 May 2026 --- 2 M}\\
      Web analytics is the process of collecting, measuring, analysing and
      reporting website data to understand and optimise web usage. Importance: it
      helps understand customer behaviour, measures campaign effectiveness,
      improves website performance and supports better marketing decisions.
      \hfill\emph{See} \S\,1.1.

\item \textbf{What is Bounce\ix{bounce} Rate in Web Analytics and how does it affect website
      performance evaluation? / How bounce rate affects website performance
      assessment.} \pyqr{CIE-II, 26 May 2026; SEE, Jun/Jul 2025 --- 2 M}\\
      Bounce rate is the percentage of visitors who leave a website after viewing
      only one page without any interaction. A high bounce rate signals poor
      content relevance, slow loading, confusing design or mismatched traffic,
      and correlates with lower conversions --- so it acts as an early diagnostic
      of both content quality and traffic quality. \hfill\emph{See} \S\,1.5.

\item \textbf{List any four important features available in Google Analytics\ix{Google Analytics} for
      website performance tracking.} \pyqr{CIE-II, 26 May 2026 --- 2 M}\\
      Traffic source analysis; bounce rate measurement; user demographics
      tracking; conversion and goal tracking. \hfill\emph{See} \S\,2.2.

\item \textbf{What is the main purpose of using Google Analytics?}
      \pyqr{CIE-II, 7 May 2025 --- 2 M}\\
      To collect, measure and report data about website users --- who they are,
      how they arrived, what they did and whether they converted --- so that
      businesses can understand behaviour, evaluate marketing performance and
      make evidence-based decisions. \hfill\emph{See} \S\,2.1.

\item \textbf{Name any two popular web analytics tools other than Google
      Analytics.} \pyqr{CIE-II, 7 May 2025 --- 2 M}\\
      Adobe Analytics and Matomo. Also acceptable: Webtrends, Mixpanel, Hotjar\ix{Hotjar},
      Microsoft Clarity, Kissmetrics, Amplitude, SimilarWeb.
      \hfill\emph{See} \S\,3.1.

\item \textbf{Define ``Dimensions'' and ``Metrics'' in Google Analytics.}
      \pyqr{SEE, Oct/Nov 2023 --- 2 M}\\
      Dimensions are attributes of the data --- qualitative descriptors such as
      city, device, browser, source and medium, or page path; they form the rows
      of a report. Metrics are quantitative measurements --- sessions, users,
      pageviews, bounce rate, revenue; they form the columns.
      \hfill\emph{See} \S\,1.4.

\item \textbf{What is meant by campaign analysis in E-mail Marketing?}
      \pyqr{CIE-II, 26 May 2026 --- 2 M}\\
      Evaluating the performance of e-mail campaigns using metrics such as open
      rate, click-through rate\ix{click-through rate}, conversion rate, bounce rate and unsubscribe
      rate, to determine effectiveness and improve future campaigns.
      \hfill\emph{See} \S\,6.1.

\item \textbf{Define Click Through Rate (CTR) in E-mail Marketing.}
      \pyqr{CIE-II, 26 May 2026 --- 2 M}\\
      CTR is the percentage of recipients who click on links in an e-mail
      compared to the total number of delivered e-mails.
      \hfill\emph{See} \S\,6.2.

\item \textbf{List any four advantages of E-mail Marketing over traditional
      marketing methods.} \pyqr{CIE-II, 26 May 2026 --- 2 M}\\
      Low cost; global reach; easy performance tracking; personalised
      communication. \hfill\emph{See} \S\,4.3.

\item \textbf{Why is segmentation important in e-mail marketing campaigns?}
      \pyqr{SEE, Jun/Jul 2025 --- 2 M}\\
      Segmentation divides the list by demographics, behaviour, engagement or
      lifecycle stage so that each group receives relevant content. It raises open
      and click rates, reduces unsubscribes and spam complaints, protects
      deliverability and enables the personalisation that drives conversion.
      \hfill\emph{See} \S\,4.5.

\item \textbf{List any two key components of a successful email marketing
      campaign.} \pyqr{SEE, Oct/Nov 2023 --- 2 M}\\
      A permission-based, well-segmented opt-in list, and a compelling subject
      line with a single clear call to action\ix{call to action}. Also acceptable: personalised
      relevant content; mobile-responsive design; measurable KPIs.
      \hfill\emph{See} \S\,4.6.

\item \textbf{After conducting an email marketing campaign for a travel agency,
      outline two specific metrics you would analyze to assess the campaign's
      performance.} \pyqr{CIE-II Quiz, 24 Jul 2024 --- 2 M}\\
      Click-through rate, to judge whether the destinations and offers presented
      were relevant enough to earn a click; and conversion rate --- bookings or
      enquiry-form completions per click --- to judge whether the campaign
      produced business value rather than merely interest.
      \hfill\emph{See} \S\,6.3.

\item \textbf{You are tasked with creating an effective email campaign to promote
      a summer sale for a fashion brand. Describe two strategies you would employ
      to increase open rates and click-through rates.}
      \pyqr{CIE-II Quiz, 24 Jul 2024 --- 2 M}\\
      A personalised, urgency-bearing, A/B-tested subject line with an optimised
      send time to raise opens; and a single dominant benefit-led call to action
      above the fold with segmented product recommendations to raise clicks.
      \hfill\emph{See} \S\,6.5.

\end{enumerate}

\section*{Part B --- Eight- and Ten-Mark Questions}
\addcontentsline{toc}{section}{Part B --- Eight- and Ten-Mark Questions}

\begin{enumerate}[leftmargin=2.2em,itemsep=1.4ex]

\item \textbf{Differentiate between traditional marketing analysis and web
      analytics-based marketing analysis with suitable examples.}
      \pyqr{CIE-II, 26 May 2026 --- 10 M}\\
      Use the comparison table --- data collection, speed, accuracy, customer
      tracking, cost, measurement, sample, attribution and feedback loop --- each
      with an example. \hfill\emph{See} \S\,1.2.

\item \textbf{Explain the role of web analytics in digital marketing. How can
      businesses benefit from analysing website data?}
      \pyqr{CIE-II, 7 May 2025 --- 10 M}\\
      The eight roles and the seven benefits. \hfill\emph{See} \S\,1.7 and
      \S\,1.8.

\item \textbf{Develop a basic strategy to set up Google Analytics for a new
      e-commerce website. What key metrics should be tracked?}
      \pyqr{CIE-II, 7 May 2025 --- 10 M}\\
      \emph{Setup strategy:} (1) create the property and data stream; (2) deploy
      the tag through Google Tag Manager\ix{Google Tag Manager} rather than hard-coding; (3) define the
      measurement plan --- what business questions must the data answer; (4)
      configure enhanced e-commerce events (view item, add to cart, begin
      checkout, purchase); (5) mark purchase and lead events as conversions and
      assign values; (6) build the checkout funnel exploration; (7) apply UTM
      tagging conventions to every campaign; (8) link the advertising platform
      and Search Console; (9) configure internal-traffic and bot filters; (10)
      set up audiences for remarketing; (11) build a dashboard; (12) validate in
      debug view and real time before trusting the data.
      \emph{Key metrics:} sessions and users by channel; new against returning;
      bounce or engagement rate; average session duration; product views,
      add-to-cart rate and cart-abandonment rate; checkout completion rate;
      conversion rate; average order value; revenue; return on advertising spend;
      and cost per acquisition by channel.
      \hfill\emph{See} \S\,2.2 and \S\,2.5.

\item \textbf{How can businesses align Web Analytics KPIs with their
      organizational goals? Illustrate with examples.}
      \pyqr{CIE-I, 16 Apr 2026 --- 10 M}\\
      The five-step method with the goal-to-KPI mapping and a SMART\ix{SMART} example.
      \hfill\emph{See} \S\,1.9.

\item \textbf{Discuss the concept of ``Acquisition channels'' in Google
      Analytics. How can business use this information to optimize their
      marketing efforts?} \pyqr{SEE, Oct/Nov 2023, Q5a --- 8 M}\\
      Define acquisition channels, list them with meaning and optimisation
      action, and explain the ways businesses act on the data --- above all,
      judging a channel by revenue rather than by traffic.
      \hfill\emph{See} \S\,2.3.

\item \textbf{Explain the significance of A/B testing\ix{A/B testing} and conversion rate
      optimization in Web Analytics. How can these practices help improve website
      performance?} \pyqr{SEE, Oct/Nov 2023, Q5b --- 8 M}\\
      Definitions, significance, what to test, the six-step process and the
      tools. \hfill\emph{See} \S\,1.10.

\item \textbf{List and explain the various reports that can be generated using
      Google Analytics.} \pyqr{SEE, Oct/Nov 2023, Q6b --- 8 M}\\
      Real-time, audience and user, acquisition, behaviour and engagement,
      conversions and monetisation, attribution, and custom reports and
      explorations. \hfill\emph{See} \S\,2.4.

\item \textbf{Describe the process of setting up and tracking goals in Google
      Analytics. How can businesses use goal tracking to measure website
      performance?} \pyqr{SEE, Oct/Nov 2023, Q10b --- 6 M}\\
      The eight-step process, and the point that goals convert traffic reports
      into ROI reports. \hfill\emph{See} \S\,2.5.

\item \textbf{Design a simple marketing plan for an e-learning platform using
      insights gathered from web analytics tools.}
      \pyqr{SEE, Jun/Jul 2025, Q5a --- 8 M}\\
      \emph{Insight gathering:} analytics shows that most traffic arrives from
      organic search\ix{organic search} on mobile, that the highest-engagement pages are free sample
      lessons, and that 70\% of drop-off occurs on the payment page; session
      recordings show users hesitating at a nine-field sign-up form.
      \emph{Plan:} (1) objective --- increase paid enrolments by 25\% in one
      quarter; (2) audience --- segment into browsers, free-trial users and
      lapsed users; (3) channels --- SEO around ``learn X online'' long-tail
      queries, video sample lessons, and a re-engagement e-mail sequence; (4)
      content --- the free sample lesson as the lead magnet, since analytics
      identified it as the strongest engagement asset; (5) conversion fixes ---
      cut the sign-up form to four fields, add trust signals and a money-back
      guarantee at the payment step, and mobile-optimise checkout; (6) e-mail ---
      trigger a nurture sequence on free-trial start and an abandoned-checkout
      reminder; (7) measurement --- track enrolment conversion rate, cost per
      enrolment, checkout completion rate and course-completion retention;
      review fortnightly. \hfill\emph{See} Chapter~1 and Chapter~5.

\item \textbf{Analyze a failed e-mail marketing campaign and identify areas where
      marketing research could have improved outcomes.}
      \pyqr{SEE, Jun/Jul 2025, Q5b --- 8 M}\\
      Diagnose against the funnel. \emph{High bounce rate} \ra{} the list was
      purchased or stale; research into list source and hygiene would have
      prevented it. \emph{Low open rate} \ra{} the subject line and sender name
      were untested; subject-line A/B research would have identified a winner.
      \emph{Low CTR despite opens} \ra{} content was irrelevant to the segment;
      audience research and segmentation would have matched offer to interest.
      \emph{Clicks but no conversions} \ra{} the landing page did not match the
      e-mail's promise; message-match testing and usability\ix{usability} research would have
      caught it. \emph{High unsubscribes and spam complaints} \ra{} expectations
      were never set at sign-up and frequency was too high; preference research
      would have fixed both.

      \smallskip
      \emph{Conclusion:} the failure was not a creative failure but a
      \textbf{research failure} --- the campaign was designed from assumption
      rather than from audience data. \hfill\emph{See} \S\,6.2 and \S\,6.3.

\item \textbf{Describe the importance of audience segmentation and
      personalization in email marketing. Provide a step-by-step guide on how to
      segment an audience effectively and personalize email content based on
      segmentation criteria.} \pyqr{CIE-II, 24 Jul 2024 --- 10 M}\\
      The eight-step guide with the fashion-retailer example, preceded by the
      importance argument. \hfill\emph{See} \S\,5.3 and \S\,4.5.

\item \textbf{Imagine you are a marketing manager for a startup selling organic
      skincare products. Develop a case study outlining the steps you would take
      to create and execute an effective e-mail campaign to launch a new product
      line. Include specific goals, target audience segmentation, content
      strategy, and metrics for success.}
      \pyqr{CIE-II, 24 Jul 2024 --- 10 M}\\
      \emph{Goals:} 3,000 pre-launch sign-ups; 25\% open rate; 4\% conversion on
      the launch send; \rupee 8 lakh revenue in the launch fortnight.
      \emph{Segmentation:} existing customers by past category (cleansers,
      serums, sun care); engaged non-buyers; customers lapsed over 180 days; and
      new sign-ups from the pre-launch teaser.
      \emph{Content strategy --- a five-e-mail sequence:} (1) a teaser to the
      full list, building anticipation with an ingredient story; (2) an
      educational e-mail on the skin concern the range solves, establishing
      authority; (3) the launch announcement with a hero product image and a
      single call to action, personalised by segment; (4) a social-proof e-mail
      carrying testimonials and before-and-after images; (5) an urgency e-mail
      for non-openers and cart abandoners with a 48-hour launch offer. Each is
      mobile-responsive, single-topic, with one dominant call to action.
      \emph{Execution:} double opt-in capture, dynamic product blocks by segment,
      send-time optimisation, and A/B tests on subject lines for e-mails 3 and 5.
      \emph{Metrics:} list growth rate, non-bounce rate, open rate,
      click-through rate, conversion rate, revenue per message, unsubscribe and
      spam-complaint rate --- reviewed after each send so the sequence is
      optimised mid-flight. \hfill\emph{See} Chapter~5 and Chapter~6.

\item \textbf{Evaluate the challenges faced in maintaining successful E-mail
      Marketing campaigns and suggest suitable solutions to improve campaign
      performance.} \pyqr{CIE-II, 26 May 2026 --- 10 M}\\
      The six challenges with paired solutions, plus deliverability, list decay,
      rendering inconsistency and privacy changes. \hfill\emph{See} \S\,6.4.

\item \textbf{Create a strategy to integrate Web Analytics and E-mail Marketing
      for improving customer engagement and business growth.}
      \pyqr{CIE-II, 26 May 2026 --- 10 M}\\
      The six-step strategy and the four benefits. \hfill\emph{See} \S\,7.2.

\item \textbf{Analyse the effectiveness of an E-mail Marketing campaign using
      KPIs such as open rate, click-through rate, unsubscribe rate, and
      conversion rate.} \pyqr{CIE-II, 26 May 2026 --- 10 M}\\
      Define each KPI, state what a high or low reading indicates, and give the
      corrective action. Anchor the answer in the six-stage measurement funnel so
      that the KPIs are shown as a connected system rather than a list.
      \hfill\emph{See} \S\,6.2 and \S\,6.3.

\item \textbf{Design an E-mail Marketing plan for a newly launched online business
      targeting college students.} \pyqr{CIE-II, 26 May 2026 --- 10 M}\\
      The official seven-point scheme answer. \hfill\emph{See} \S\,5.2.

\item \textbf{Evaluate the effectiveness of responsive email design in
      maintaining a competitive edge in digital marketing.}
      \pyqr{SEE, Jun/Jul 2025, Q6a --- 8 M}\\
      The case for it, how it is achieved, and the critical view that it is now
      table stakes rather than a differentiator. \hfill\emph{See} \S\,7.4.

\item \textbf{Critically evaluate the contribution of Google Tag Manager in
      optimizing marketing channels for global companies.}
      \pyqr{SEE, Jun/Jul 2025, Q6b --- 8 M}\\
      The six contributions and the four critical limitations.
      \hfill\emph{See} \S\,2.6.

\item \textbf{Demonstrate how tools like Google Analytics and Hotjar can be used
      to develop insights into user behavior and improve marketing decisions.}
      \pyqr{SEE, Jun/Jul 2025, Q7a --- 8 M}\\
      Analytics gives the \emph{what} --- which channels deliver traffic, which
      landing pages bounce, where users drop out of checkout, which segments
      convert, and what each channel costs per acquisition. Behavioural tools
      give the \emph{why} --- heatmaps show which elements attract attention,
      scroll maps show whether users ever reach the call to action, click maps
      expose rage clicks on non-clickable elements, session recordings show
      hesitation in real time, and on-page surveys capture the reason for
      abandonment in the user's own words. Then give the combined workflow and
      the decisions it improves. \hfill\emph{See} \S\,3.2.

\end{enumerate}

\section*{How Unit III Is Examined}
\addcontentsline{toc}{section}{How Unit III Is Examined}

\begin{keybox}[The Pattern Across Papers]
In the Semester End Examination, Unit~III is Questions~5 and~6 with internal
choice --- 16 marks in two eight-mark sub-divisions --- plus two or three Part~A
objective questions. The second Continuous Internal Evaluation test draws heavily
on it.

\smallskip
The highest-probability topics are the \textbf{e-mail measurement funnel and KPI
analysis}, \textbf{designing an e-mail marketing plan} for a given scenario,
\textbf{segmentation and personalisation}, \textbf{Google Analytics reports, goals
and acquisition channels}, and \textbf{traditional against web-analytics-based
analysis}.

\smallskip
Note that almost every Unit~III question is \emph{applied} --- you will be given a
business and asked to build or diagnose something --- so always answer with a
named framework (the seven-node plan, the six-stage funnel, the eight-step goal
setup) rather than general commentary.
\end{keybox}

\chapter{Glossary of Key Terms}

\begin{mmlongtable}{Key terms for rapid revision}{tab:glossary3}%
{@{}P{0.26\textwidth}P{0.67\textwidth}@{}}
\rowcolor{ocre}\thd{Term} & \thd{One-line meaning}\\
\midrule
\endfirsthead
\rowcolor{ocre}\thd{Term} & \thd{One-line meaning}\\
\midrule
\endhead
\rlab{A/B testing} & Splitting traffic between two versions to find the
statistically better one\\
\rowcolor{tablealt}
\rlab{Acquisition channel} & The grouping describing where a session came from\\
\rlab{Auto-responder} & An automated sequence released to new subscribers\\
\rowcolor{tablealt}
\rlab{Bounce rate (web)} & Percentage of visitors who leave after viewing one
page without interaction\\
\rlab{Bounce rate (e-mail)} & Percentage of sent e-mails that were not
delivered\\
\rowcolor{tablealt}
\rlab{Buzz} & A combined popularity measure across social mentions\\
\rlab{CAN-SPAM} & US anti-spam law; fines up to \$16,000 per violation\\
\rowcolor{tablealt}
\rlab{Campaign analysis} & Evaluating e-mail performance on opens, clicks,
conversions, bounces and unsubscribes\\
\rlab{CRO} & Conversion rate optimisation --- systematically raising the
percentage who convert\\
\rowcolor{tablealt}
\rlab{CTR} & Percentage of recipients, or of opens, who clicked a link\\
\rlab{Dimension} & An attribute of the data --- city, device, source; forms the
rows of a report\\
\rowcolor{tablealt}
\rlab{Double opt-in} & Consent confirmed through a verification e-mail\\
\rlab{Drip campaign} & A scheduled nurture sequence delivered over time\\
\rowcolor{tablealt}
\rlab{E-mail marketing} & Using e-mail to deliver commercial and relationship
messages to a permission-based list\\
\rlab{Exit rate} & Percentage of sessions that ended on a given page\\
\rowcolor{tablealt}
\rlab{Funnel visualisation} & A report showing drop-off at each step of a defined
path\\
\rlab{Goal / conversion} & A defined, valuable user action being tracked\\
\rowcolor{tablealt}
\rlab{Google Analytics} & The platform that collects, measures and reports
website and app user data\\
\rlab{Google Tag Manager} & A container for deploying and versioning tracking
tags without code changes\\
\rowcolor{tablealt}
\rlab{List decay} & The natural annual loss of list quality, roughly 20--25\% per
year\\
\rlab{Metric} & A quantitative measurement --- sessions, revenue; forms the
columns of a report\\
\rowcolor{tablealt}
\rlab{Non-bounce rate} & $100\%$ minus bounced over total sent --- the delivery
rate\\
\rlab{Open rate} & Percentage of delivered e-mails opened\\
\rowcolor{tablealt}
\rlab{Opt-in} & Permission willingly given to receive e-mail\\
\rlab{Pageviews} & Total pages loaded; unique pageviews counts a page once per
session\\
\rowcolor{tablealt}
\rlab{Revenue per message} & Total revenue divided by the number of sent
e-mails\\
\rlab{Segment} & An isolated subset of users defined by shared attributes or
behaviour\\
\rowcolor{tablealt}
\rlab{Session} & A group of user interactions within a given time frame\\
\rlab{Share of voice} & Conversations about you against conversations about
competitors\\
\rowcolor{tablealt}
\rlab{Transactional e-mail} & A user-triggered informational e-mail --- receipt,
one-time password, shipping notice\\
\rlab{Trigger} & An automated e-mail fired by a specific user action\\
\rowcolor{tablealt}
\rlab{Unsubscribe rate} & Percentage of recipients who opted out\\
\rlab{Web analytics} & Collecting, measuring, analysing and reporting website
data to understand and optimise web usage\\
\end{mmlongtable}

\chapter{Framework Quick-Reference Sheet}

Everything in this volume that is a numbered list, in one place, for the night
before the examination.

\begin{mmlongtable}{Framework quick reference}{tab:quickref3}%
{@{}P{0.28\textwidth}P{0.65\textwidth}@{}}
\rowcolor{ocre}\thd{Framework} & \thd{The sequence}\\
\midrule
\endfirsthead
\rowcolor{ocre}\thd{Framework} & \thd{The sequence}\\
\midrule
\endhead
\rlab{Three metric families} & Volume (are we seen?), quality (are we relevant?),
action (are they engaging?)\\
\rowcolor{tablealt}
\rlab{Traditional against web analytics} & Data collection, speed, accuracy,
customer tracking, cost, measurement, sample, attribution, feedback loop\\
\rlab{Bounce against exit} & Bounce is a single-page no-interaction session; exit
is the last page of any session\\
\rowcolor{tablealt}
\rlab{Dimensions and metrics} & Attributes form the rows; numbers form the
columns\\
\rlab{Eight roles of web analytics} & Effectiveness, behaviour, conversion,
budget, content, technical, segmentation, forecasting\\
\rowcolor{tablealt}
\rlab{Seven benefits} & Conversion, acquisition cost, experience, problem
detection, content decisions, ROI proof, competitive advantage\\
\rlab{KPI alignment (five steps)} & Identify goals \ra{} map to KPIs \ra{} make
SMART \ra{} instrument analytics \ra{} optimise continuously\\
\rowcolor{tablealt}
\rlab{CRO process} & Find the weakest step \ra{} hypothesise \ra{} build the
variant \ra{} split traffic \ra{} run to significance \ra{} implement \ra{}
repeat\\
\rlab{Four GA features to name} & Traffic source analysis, bounce rate,
demographics, conversion and goal tracking\\
\rowcolor{tablealt}
\rlab{Nine acquisition channels} & Organic search, paid search, direct, referral,
organic social, paid social, e-mail, display, affiliates\\
\rlab{Seven report families} & Real-time, audience, acquisition, behaviour,
conversions, attribution, custom and explorations\\
\rowcolor{tablealt}
\rlab{Goal setup (eight steps)} & Define the conversion \ra{} choose the type
\ra{} implement \ra{} assign value \ra{} configure the funnel \ra{} verify \ra{}
track \ra{} act\\
\rlab{Tag Manager --- for and against} & Faster deployment, standard data layer,
consent, page weight, experimentation, rollback \emph{against} tag sprawl,
duplicate firing, performance, security\\
\rowcolor{tablealt}
\rlab{What against why} & Analytics locates the drop-off; heatmaps, recordings and
surveys explain it\\
\rlab{E-mail 1.0 against 2.0} & One-way, static, batch, untracked \emph{against}
triggered, segmented, automated, measured\\
\rowcolor{tablealt}
\rlab{Four e-mail advantages} & Low cost, global reach, easy tracking,
personalised communication\\
\rlab{Opt-in against spam} & Acquisition method, expectation setting, brand
impact, list quality\\
\rowcolor{tablealt}
\rlab{Three e-mail rules} & One topic, simple design complementing the message,
one clear call to action\\
\rlab{Nine e-mail types} & Revenue: promotional, new inventory, reorder.
Relationships: newsletter, welcome, educational, product advice. Intelligence:
testimonials, surveys\\
\rowcolor{tablealt}
\rlab{Transactional against promotional} & Trigger, purpose, consent, volume,
open rate (40--60\% against 15--25\%), examples, timing\\
\rlab{Seven-node e-mail plan} & Software \ra{} list \ra{} capture forms \ra{}
goals \ra{} auto-responders \ra{} triggers \ra{} metrics \ra{} back to software\\
\rowcolor{tablealt}
\rlab{Segmentation (eight steps)} & Collect data \ra{} choose criteria \ra{}
build dynamic segments \ra{} design per segment \ra{} personalise within \ra{}
automate triggers \ra{} test within \ra{} measure per segment\\
\rlab{Six-stage measurement funnel} & Sent \ra{} delivered \ra{} opened \ra{}
clicked \ra{} converted \ra{} revenue per message\\
\rowcolor{tablealt}
\rlab{Six e-mail challenges} & Low opens, spam filtering, unsubscribes, poor
design, no personalisation, inactive subscribers\\
\rlab{Offline tracking} & Point-of-sale e-mail capture, dedicated phone numbers,
e-mail-only entrance codes; promo codes, UTM tags, QR coupons\\
\rowcolor{tablealt}
\rlab{Analytics--e-mail integration} & Collect \ra{} segment \ra{} personalise
\ra{} track \ra{} retarget \ra{} optimise\\
\rlab{Four pillars} & The fuel (opt-in list), the payload (content arsenal), the
engine (triggers and tech), the dashboard (metrics)\\
\end{mmlongtable}
